\subsection{环分解的应用}

引理 2. 设 A 是一个环，m 是 A 的理想，且有序对 (A, m) 满足 Hensel 条件。设 B 是商环 A/m，\(\pi: A \to B\) 是典范同态。对于 B 的每个幂等元 c，存在唯一的 A 的幂等元 e，使得 \(\pi(e) = c\)。

取 \(a\) 使得 \(\pi(a) = c\)；将第 5 节定理 2 的推论 1 应用于多项式 \(f(\mathbf{X}) = \mathbf{X}^2 - \mathbf{X}\)，它属于 \(\mathbf{A}[\mathbf{X}]\)，以及元素 \(a \in \mathbf{A}\)。则 \(f'(a) = 2a - 1\)；由于 \(\pi(2a - 1) = 2c - 1\)，且 \((2c - 1)^2 = 1\) 在 \(\mathbf{B}\) 中成立，\(2c - 1\) 在 \(\mathbf{B}\) 中可逆，因而 \(2a - 1\) 在 \(\mathbf{A}\) 中可逆（§ 2，第 13 节，引理 3）。由于 \(f(a) \in \mathfrak{m}\)，第 5 节定理 2 的推论 1 立即给出 \(e\) 的存在性和唯一性。

命题 8. 设 A 是一个环，m 是 A 的理想，且有序对 (A, m) 满足 Hensel 条件。设 B 是商环 A/m，π: A → B 是典范同态。若 B 是理想有限族 \((\mathfrak{b}_{i})_{i\in I}\) 的直分解，则存在唯一的 A 的理想族 \((\mathfrak{a}_{i})_{i\in I}\)，使得 \(\pi(\mathfrak{a}_{i}) = \mathfrak{b}_{i}\) 对所有 \(i \in I\) 成立，且 A 是理想族 \((\mathfrak{a}_{i})\) 的直分解。

令 \(1 = \sum_{i} c_{i}\)，其中 \(c_{i} \in \mathfrak{b}_{i}\)（对所有 \(i\)）；\(c_{i}\) 是 \(\mathbf{B}\) 中的幂等元，满足 \(c_{i}c_{j} = 0\)（当 \(i \neq j\) 时）。由引理 2，存在幂等元 \(e_{i}\)，它们属于 \(\mathbf{A} (i \in \mathbf{I})\)，使得 \(\pi(e_{i}) = c_{i}\)（对所有 \(i\)）；由于 \(e_{i}e_{j}\) 是一个幂等元，满足

\[
\pi (e _ {i} e _ {j}) = c _ {i} c _ {j} = 0
\]

对于 \(i \neq j\)，有 \(e_i e_j = 0\)，这对 \(i \neq j\) 成立（引理 2）；由于 \(1 - \sum_{i} e_i\) 是满足 \(\pi\left(1 - \sum_{i} e_i\right) = 1 - \sum_{i} c_i = 0\) 的幂等元，同样有 \(1 = \sum_{i} e_i\)。于是 \(A\) 是理想 \(a_i = e_i A\) 的直分解，且 \(\pi(a_i) = \pi(e_i) B = b_i\)。

还需证明这种分解的唯一性。现在，设 A 是另一理想族 \((\mathfrak{a}_i')_{i\in I}\) 的直分解，且 \(\pi (\mathfrak{a}_i^{\prime}) = \mathfrak{b}_i\) 对所有 \(i\) 成立；于是 \(1 = \sum_{i}e_{i}^{\prime}\)，其中 \(e_i^\prime \in \mathfrak{a}_i^\prime\)，从而在 B 中 \(1 = \sum_{i}\pi (e_i^{\prime})\)，其中 \(\pi (e_i^{\prime})\in \mathfrak{b}_i\)，这蕴含 \(\pi (e_i^{\prime}) = c_i\)；由于 \(e_i^\prime\) 和 \(e_i\) 都是幂等元，必有 \(e_i^\prime = e_i\)（引理 2），证明完毕。

注。命题 8 再次给出了 A 的半局部环结构，其中 A 关于 r-进拓扑是 Hausdorff 且完备的（r 是 A 的 Jacobson 根）；这一结构此前已作为 § 2 第 13 节命题 19 的推论得到。

